\documentclass[11pt,a4paper]{article}

\usepackage{thga-db}
\usepackage{caption}   % enables \captionof inside minipages
\usepackage{tikz}
\usetikzlibrary{positioning}

% ----- Placeholder figure macro ----
\newcommand{\placeholder}[3]{%
  \begin{tikzpicture}
    \node[draw=THGABlau!45, fill=THGABlau!5, very thick, rounded corners=2pt,
          minimum width=#1, minimum height=#2, align=center, inner sep=6pt,
          text=THGABlau!75, font=\small\itshape] {#3};
  \end{tikzpicture}%
}

% ----- Metadata --------------------------------------------
\renewcommand{\thgaKurs}{Databases}
\renewcommand{\thgaDatum}{Summer Term 2026}
\newcommand{\authorName}{Baha Eddine Nabhan}
\newcommand{\projectTitle}{LabTrack: University Laboratory Equipment \& Loan Management System}

\begin{document}
\thispagestyle{empty}
\thgaTitel{Term Project – Proposal}{\projectTitle}

\begin{infobox}
\textbf{Author:} \authorName \\
\textbf{Module:} Introduction to Database Management Systems · THGA Bochum \\
\textbf{Submitted to:} Stephan Bökelmann · \href{mailto:sboekelmann@ep1.rub.de}{sboekelmann@ep1.rub.de} \\
\textbf{Target size:} about 6 pages
\end{infobox}

\begin{infobox}
\textbf{Every figure must be explained in the running text.} A figure never
stands on its own: refer to it by number and say what it shows -- e.g.
``Figure~\ref{fig:er} shows the ER model of \dots''. The placeholder boxes below
mark where your own images go; replace each with \texttt{\textbackslash
includegraphics} and keep the caption.
\end{infobox}

% ============================================================
\section{Application domain and motivation}
% ============================================================

In university research labs, tracking hardware inventory (such as oscilloscopes, microcontrollers, and sensors) across various courses and student projects often relies on manual spreadsheets. This manual approach leads to untracked equipment loss, double-booking conflicts, and a lack of visibility into device maintenance states.

\textbf{LabTrack} is a central laboratory management system designed for lab administrators and university students.
\begin{itemize}
    \item \textbf{Scenario 1 (Student Loan Request):} A student logs into the desktop application to check the real-time availability of measurement equipment for an upcoming practical exam. They submit a loan request specifying the start and target return dates.
    \item \textbf{Scenario 2 (Admin Management \& Auditing):} A lab manager uses protected admin endpoints to approve equipment loans, log hardware condition maintenance, and inspect high-usage assets via aggregated borrowing metrics.
\end{itemize}

% ============================================================
\section{Data model -- ER sketch}
% ============================================================

\begin{figure}[ht]
  \centering
  % Replace with: \includegraphics[width=0.85\linewidth]{er-sketch.jpg}
  \placeholder{0.85\linewidth}{4.5cm}{[ ER Model Sketch ]\\ Entity Types: Student, Equipment, Loan, MaintenanceLog\\ N:M Relationship: Student <-> Equipment via Loan}
  \caption{ER sketch of the LabTrack data model showing entities, key attributes, cardinalities, and the N:M loan relationship.}
  \label{fig:er}
\end{figure}

Figure~\ref{fig:er} illustrates the conceptual ER model for the LabTrack domain. The central entity types are \texttt{Student}, \texttt{Equipment}, \texttt{Loan}, and \texttt{MaintenanceLog}. A student can request multiple loans over time, while an item of equipment can be borrowed repeatedly across its lifecycle, establishing a classic $N:M$ relationship between \texttt{Student} and \texttt{Equipment}. This $N:M$ link is explicit through the junction entity \texttt{Loan}, which tracks key transactional attributes such as \texttt{start\_date}, \texttt{due\_date}, and \texttt{status}. Additionally, \texttt{Equipment} has a $1:N$ relationship with \texttt{MaintenanceLog} to track service records.

% ============================================================
\section{From model to schema}
% ============================================================

The relational schema is derived directly from the conceptual ER diagram:

\begin{itemize}
  \item \texttt{students}(\underline{student\_id}, name, email, matriculation\_no)
  \item \texttt{categories}(\underline{category\_id}, category\_name, location\_rack)
  \item \texttt{equipment}(\underline{equipment\_id}, serial\_no, model\_name, status, \emph{category\_id})
  \item \texttt{loans}(\underline{loan\_id}, \emph{student\_id}, \emph{equipment\_id}, issue\_date, due\_date, return\_date)
  \item \texttt{maintenance\_logs}(\underline{log\_id}, \emph{equipment\_id}, log\_date, issue\_description, cost)
\end{itemize}

\textbf{Normalisation.} The schema satisfies the Third Normal Form (3NF). Every non-key attribute depends strictly on its primary key (1NF/2NF), and there are no transitive functional dependencies between non-key attributes (e.g., equipment category locations are isolated in the \texttt{categories} table to prevent duplicate storage).

% ============================================================
\section{API design}
% ============================================================

The FastAPI backend exposes structured endpoints separating public read access from privileged operations requiring an \texttt{X-API-Key}:

\renewcommand{\arraystretch}{1.2}
\begin{tabularx}{\linewidth}{@{}p{2.2cm}p{4.8cm}X@{}}
\toprule
\textbf{Method} & \textbf{Path} & \textbf{Purpose} \\
\midrule
\texttt{GET}  & \texttt{/api/v1/equipment} & Reads active equipment list with category details (\texttt{JOIN}). \\
\texttt{GET}  & \texttt{/api/v1/analytics/usage} & Computes total loan count and uptime ratio per category (\texttt{JOIN} + aggregation). \\
\texttt{POST} & \texttt{/api/v1/loans} & Creates a new equipment loan record (requires \texttt{X-API-Key}). \\
\texttt{POST} & \texttt{/api/v1/maintenance} & Adds a maintenance log entry (requires \texttt{X-API-Key}). \\
\bottomrule
\end{tabularx}

% ============================================================
\section{Frontend prototype (hand sketches)}
% ============================================================

The frontend client is implemented using Python and \texttt{tkinter}. It is packaged and distributed as a native \texttt{.deb} installer for Linux desktop environments. On startup, the application prompts the user for authentication credentials and the target backend endpoint url.

\begin{figure}[ht]
\centering
\begin{minipage}[t]{0.47\linewidth}\centering
  \placeholder{\linewidth}{3.5cm}{[ Hand Sketch 1 ]\\ Connection Dialog}
  \captionof{figure}{Startup connection dialog collecting the API URL and X-API-Key.}
  \label{fig:fe-connect}
\end{minipage}\hfill
\begin{minipage}[t]{0.47\linewidth}\centering
  \placeholder{\linewidth}{3.5cm}{[ Hand Sketch 2 ]\\ Equipment List}
  \captionof{figure}{Main screen: searches inventory via \texttt{GET /api/v1/equipment}.}
  \label{fig:fe-main}
\end{minipage}

\vspace{1em}

\begin{minipage}[t]{0.47\linewidth}\centering
  \placeholder{\linewidth}{3.5cm}{[ Hand Sketch 3 ]\\ Loan Creation Form}
  \captionof{figure}{Detail form: creates loan via \texttt{POST /api/v1/loans}.}
  \label{fig:fe-detail}
\end{minipage}\hfill
\begin{minipage}[t]{0.47\linewidth}\centering
  \placeholder{\linewidth}{3.5cm}{[ Hand Sketch 4 ]\\ Analytics Dashboard}
  \captionof{figure}{Analytics view: calls \texttt{GET /api/v1/analytics/usage}.}
  \label{fig:fe-second}
\end{minipage}
\end{figure}

Figures~\ref{fig:fe-connect}--\ref{fig:fe-second} illustrate the four wireframe screens of the GUI application. Figure~\ref{fig:fe-connect} shows the connection modal displayed upon startup to capture backend connectivity details. Figure~\ref{fig:fe-main} displays the primary dashboard, presenting available equipment fetched via \texttt{GET /api/v1/equipment}. Figure~\ref{fig:fe-detail} displays the submission dialog for reserving equipment via \texttt{POST /api/v1/loans}. Finally, Figure~\ref{fig:fe-second} depicts the administration panel visualizing equipment loan frequency built using aggregate data.

% ============================================================
\section{Operation with Docker Compose}
% ============================================================

The application infrastructure runs using \texttt{docker-compose.yml} managing two main services:
\begin{enumerate}
    \item \texttt{db}: A PostgreSQL container running image \texttt{postgres:16-alpine} with persistent volume mounts.
    \item \texttt{api}: A Python 3.11 container hosting the FastAPI server via \texttt{uvicorn}.
\end{enumerate}
Database credentials, port assignments, and API secret keys are stored in a local \texttt{.env} file (excluded from version control). For production grading on the lecture server, deployment is initiated via \texttt{docker compose up -d --build}.

% ============================================================
\section{Scope, milestones and risks}
% ============================================================

\textbf{Scope.} 5 relational tables, 4 FastAPI endpoints, 1 \texttt{tkinter} desktop GUI client packaged as a \texttt{.deb} file.

\textbf{Milestones and schedule.}

\renewcommand{\arraystretch}{1.2}
\begin{tabularx}{\linewidth}{@{}p{1.4cm}Xp{1.8cm}@{}}
\toprule
\textbf{Week} & \textbf{Milestone} & \textbf{Est. hours} \\
\midrule
1   & PostgreSQL relational schema setup, DDL scripts, and mock seed data. & 8\,h \\
2--3 & FastAPI backend implementation with ORM/SQL, joins, aggregation, and authentication. & 12\,h \\
4--5 & \texttt{tkinter} GUI client development, API interaction, and \texttt{.deb} package assembly. & 10\,h \\
6   & Automated integration testing, documentation, and video demonstration recording. & 10\,h \\
\bottomrule
\end{tabularx}

\textbf{Risks.} A primary risk is hardware status desynchronization caused by simultaneous loan submissions. This is mitigated at the database layer using explicit transactions with strict isolation levels or check constraints.

% ============================================================
\section{Acceptance criteria and testing}
% ============================================================

\textbf{Acceptance criteria.}
\begin{itemize}[topsep=2pt]
  \item Attempting to issue a loan for an item marked \texttt{'MAINTENANCE'} is rejected by the API with a \texttt{400 Bad Request}.
  \item \texttt{GET /api/v1/analytics/usage} correctly calculates aggregated usage counts joining \texttt{equipment} and \texttt{loans}.
  \item Write requests (\texttt{POST /api/v1/loans}) missing a valid \texttt{X-API-Key} header fail immediately with HTTP \texttt{401 Unauthorized}.
\end{itemize}

\textbf{Test plan.}
\begin{itemize}[topsep=2pt]
  \item \textbf{Database:} Verify primary key, foreign key CASCADE integrity, and \texttt{CHECK} constraints using SQL integration scripts.
  \item \textbf{API:} Automated endpoint test suite built with \texttt{pytest} and FastAPI's \texttt{TestClient} targeting an isolated test database.
  \item \textbf{Business rule:} Unit tests asserting that double-booking or illegal state transitions trigger transactional rollbacks.
\end{itemize}

\vspace{1em}
\begin{infobox}
\textbf{Effort budget:} the whole term project (system, documentation and video)
should fit into roughly \textbf{40\,h of work} over at most \textbf{2 months}.
\end{infobox}

\end{document}
